\exercisesection{赋范代数}

\begin{enumerate}
\def\labelenumi{\arabic{enumi})}
\tightlist
\item
  设 (A, e) 为一个复幺代数。
\end{enumerate}

\begin{enumerate}
\def\labelenumi{\alph{enumi})}
\tightlist
\item
  对于线性型 \(\lambda \in \mathrm{A}^*\)，当 \(\lambda(e) = 1\) 时，下列条件等价：
\end{enumerate}

\begin{enumerate}
\def\labelenumi{(\roman{enumi})}
\item
  对每个 \(x \in \operatorname{Ker}(\lambda)\)，有 \(x^2 \in \operatorname{Ker}(\lambda)\)；
\item
  对每个 \(x \in \mathbb{A}\)，有 \(\lambda(x^{2}) = \lambda(x)^{2}\)；
\item
  \(\lambda\) 的核是 A 的一个右理想；
\item
  线性型 \(\lambda\) 是从 A 到 C 的保单位元态射。
\end{enumerate}

（注意 \(\lambda(x - \lambda(x)e) = 0\) 对所有 \(x\in A\) 成立。）

\begin{enumerate}
\def\labelenumi{\alph{enumi})}
\setcounter{enumi}{1}
\item
  设 \(\lambda \in \mathrm{A}^*\) 为满足 \(\lambda(1) = 1\) 的线性型，且使得对每个 \(x \in \mathrm{A}\)，有 \(\lambda(x) \in \mathrm{Sp}_{\mathrm{A}}(x)\)。设 \(x \in \mathrm{Ker}(\lambda)\) 满足 \(\lambda(x^{2}) \neq 0\)，则 \(x\) 的谱在 \(\mathbf{C}\) 中无界。（取整数 \(n \geqslant 2\)，考虑多项式 \(p \in \mathbf{C}[\mathrm{X}]\)，使 \(p(z) = \lambda((ze - x)^{n})\) 对所有 \(z \in \mathbf{C}\) 成立；注意它的根都属于 \(x\) 的谱，并用 \(|\lambda(x^{2})|\) 给出它们的欧氏范数的下界。）
\item
  假设 A 是 Banach 代数。为使 A 上的线性型 \(\lambda\) 是从 A 到 C 的保单位元态射，充要条件是 \(\lambda(e)=1\)，且 \(\lambda(x)\in\mathrm{Sp}_{\mathrm{A}}(x)\) 对所有 \(x\in A\) 成立。
\end{enumerate}

\begin{enumerate}
\def\labelenumi{\arabic{enumi})}
\setcounter{enumi}{1}
\tightlist
\item
  设 \(n \geqslant 0\) 为整数。设 \(A_n\) 为由函数 \(f: [0,1] \to \mathbf{C}\)（它在 \([0,1]\) 上具有直到 \(n\) 阶的连续导数）所组成的交换复 Banach 代数，赋以范数
\end{enumerate}

\[
\| f \| = \sum_ {k = 0} ^ {n} \frac {1}{k !} \sup _ {0 \leqslant t \leqslant 1} | f ^ {(k)} (t) |
\]

（I，第 18 页，例 4）。证明 \(A_{n}\) 同构于 \(A_{m}\) 当且仅当 n = m。（不过，根据 I，第 144 页，例 1，这些 Banach 代数都有相同的特征标空间。）

\begin{enumerate}
\def\labelenumi{\arabic{enumi})}
\setcounter{enumi}{2}
\tightlist
\item
  设 G 为等于其导出子群的群（A，I，第 67 页）。对每个 \(g \in G\)，用 \(\ell(g)\) 表示满足下列条件的最小整数 \(k \geqslant 0\)：存在数对 \((a_{1}, b_{1}), \ldots, (a_{k}, b_{k})\)（取自 \(G \times G\)），使 \(g = [a_{1}, b_{1}] \cdots [a_{k}, b_{k}]\)。
\end{enumerate}

\begin{enumerate}
\def\labelenumi{\alph{enumi})}
\item
  映射 \(\ell\) 是 G 上的一个度量。
\item
  对每个 \(g \in \mathbf{G}\)，极限
\end{enumerate}

\[
\operatorname{lsc} (g) = \lim _ {n \to + \infty} \frac {\ell (g ^ {n})}{n}
\]

存在，称为 g 的稳定换位子长度。

\begin{enumerate}
\def\labelenumi{\alph{enumi})}
\setcounter{enumi}{2}
\tightlist
\item
  设 S 为一个集合，\(\mathrm{G} = \mathrm{D}(\mathrm{F}(\mathrm{S}))\) 为由 S 生成的自由群的导出子群。群 G 等于其导出子群，且对每个 \(g \in G\)，\(g \neq e\)，有 \(\operatorname{lsc}(g) \geqslant 1/2\)。存在 \(g \in G\) 使 \(\operatorname{lsc}(g) = 1/2\)。
\end{enumerate}

\begin{enumerate}
\def\labelenumi{\arabic{enumi})}
\setcounter{enumi}{3}
\tightlist
\item
  设 \(n \geqslant 1\) 为整数。用 \(S(n)\) 表示满足下列条件的最大整数 \(k \geqslant 0\)：对每个集合 \(X\)（由 \(n\) 个实数组成），存在子集 \(Y \subset X\)（基数为 \(k\)），使方程 \(x + y = z\) 在 \(Y\) 中无解。于是 \(S(n) \leqslant n\)。
\end{enumerate}

\begin{enumerate}
\def\labelenumi{\alph{enumi})}
\tightlist
\item
  证明极限
\end{enumerate}

\[
\sigma = \lim _ {n \to + \infty} \frac {\mathrm{S} (n)}{n}
\]

存在。

\begin{enumerate}
\def\labelenumi{\alph{enumi})}
\setcounter{enumi}{1}
\item
  证明 \(\sigma \leqslant 2 / 5\)。
\item
  证明 \(\sigma \geqslant 1 / 3\)。
\end{enumerate}

（可以证明 \(\sigma = 1/3\)，参见 S. EBERHARD, B. GREEN, F. MANNERS, Sets of integers with no large sum-free subset, Annals of Mathematics 180 (2014), 621--652。）

\begin{enumerate}
\def\labelenumi{\arabic{enumi})}
\setcounter{enumi}{4}
\tightlist
\item
  设 A 为交换环，\(n \geqslant 1\) 为整数。对于系数在 A 中的矩阵 \(M = (m_{i,j})\)（其型为 \((n, n)\)），我们把 M 的积和式称为元素
\end{enumerate}

\[
\mathrm{per(M)} = \sum_ {\sigma \in \mathrm{S} _ {n}} m _ {1, \sigma (1)} \dots m _ {n, \sigma (n)}
\]

该元素属于 A。

设 \(k \geqslant 1\) 为整数。用 \(X_{n,k}\) 表示由实系数的、型为 \((n, n)\) 且每列恰含 k 个等于 1 的系数（其余系数全为零）的矩阵 M 所成的集合。

证明极限

\[
\lim _ {n \to + \infty} \left(\inf _ {\mathrm{M} \in \mathrm{X} _ {n, k}} \operatorname{per} (\mathrm{M})\right) ^ {1 / n}
\]

存在。

\begin{enumerate}
\def\labelenumi{\arabic{enumi})}
\setcounter{enumi}{5}
\tightlist
\item
  设 \(k \geqslant 1\) 为整数。对每个整数 \(N \geqslant 1\)，用 \(S_k(N)\) 表示不包含 \(A \subset \{1, \ldots, N\}\) 个成等差数列的元素的集合 \(k\) 的最大基数；也就是说，对任意 \(n_0 \geqslant 0\)、\(h \geqslant 1\)、\(\ell \geqslant 1\)，若 \(\{n_0 + h, n_0 + 2h, \ldots, n_0 + \ell h\} \subset A\)，则 \(\ell < k\)。
\end{enumerate}

\begin{enumerate}
\def\labelenumi{\alph{enumi})}
\tightlist
\item
  证明极限
\end{enumerate}

\[
s _ {k} = \lim _ {\mathrm{N} \to + \infty} \frac {\mathrm{S} _ {k} (\mathrm{N})}{\mathrm{N}}
\]

存在。

\begin{enumerate}
\def\labelenumi{\alph{enumi})}
\setcounter{enumi}{1}
\tightlist
\item
  证明 \(s_{1}=s_{2}=0\)。可以证明（“Szemerédi 定理”，参见 E. SZEMERÉDI, On sets of integers containing no k elements in arithmetic progression, Acta Arithmetica XXVII, 1975, 199--245）：对所有 \(s_{k}=0\) 有 \(k\geqslant1\)；至于 k=3 的情形，参见 II，第 270 页，习题 23。
\end{enumerate}

\begin{enumerate}
\def\labelenumi{\arabic{enumi})}
\setcounter{enumi}{6}
\tightlist
\item
  设 A 为赋范代数。
\end{enumerate}

\begin{enumerate}
\def\labelenumi{\alph{enumi})}
\tightlist
\item
  设 S 为 A 的非空有界子集。证明极限
\end{enumerate}

\[
\varrho (\mathrm{S}) = \lim _ {n \to + \infty} \sup _ {x \in \mathrm{S} ^ {n}} \| x \| ^ {1 / n}
\]

存在，并且当 S 仅由单个元素 x 组成时，有 \(\varrho(\mathrm{S}) = \varrho(x)\)。我们称 \(\varrho(\mathrm{S})\) 为子集 S 的谱半径。

\begin{enumerate}
\def\labelenumi{\alph{enumi})}
\setcounter{enumi}{1}
\tightlist
\item
  设 N 为 A 上与 A 的范数等价、且使赋以 p 之后的 A 仍为赋范代数的范数 p 全体所成的集合。设 X 为 A 的一个子集。为存在 \(p \in N\) 使得对所有 \(p(x) \leqslant 1\) 都有 \(x \in X\)，充要条件是存在实数 \(c \geqslant 0\) 使得
\end{enumerate}

\[
\sup _ {n \geqslant 1} \sup _ {x \in X ^ {n}} \| x \| \leqslant c.
\]

\begin{enumerate}
\def\labelenumi{\alph{enumi})}
\setcounter{enumi}{2}
\tightlist
\item
  证明
\end{enumerate}

\[
\varrho (\mathrm{S}) = \inf _ {p \in \mathrm{N}} \sup _ {x \in \mathrm{S}} p (x).
\]

\begin{enumerate}
\def\labelenumi{\arabic{enumi})}
\setcounter{enumi}{7}
\item
  设 A 为 \(\mathbf{R}\) 上在无穷远处趋于 0 的复值连续函数所成的代数，赋以范数 \(\|f\| = \sup_{t\in\mathbf{R}} |f(t)|\)。于是 \(\widetilde{\mathrm{A}}\) 就恒同于 \(\mathbf{R}\) 上在无穷远处趋于有限极限的复值连续函数所成的代数。在 \(\widetilde{\mathrm{A}}\) 上，范数 \(\|g\| = \sup_{t\in\mathbf{R}} |g(t)|\) 与 I，第 15 页，第 1 节中定义的范数不同。
\item
  设 A 为赋范代数，\(b = \inf_{x \neq 0} \|x^{2}\|/\|x\|^{2}\)。则
\end{enumerate}

\[
b \leqslant \inf _ {x \neq 0} \frac {\varrho (x)}{\| x \|} \leqslant \sqrt {b}.
\]

\begin{enumerate}
\def\labelenumi{\arabic{enumi})}
\setcounter{enumi}{9}
\item
  设 A 为赋范代数，\(x, y \in A\)。我们有 \(\varrho(xy) = \varrho(yx)\)。
\item
  设 H 为可数型 Hilbert 空间，\((f_{n})_{n\geqslant1}\) 为 H 的一个标准正交基。设 A 为赋范代数 \(\mathcal{L}(\mathrm{H})\)。设 \(x\in A\) 为由 \(x(f_{i})=2^{-i}f_{i+1}\) 所定义的元素。证明 x 是拟幂零元，但不是幂零元。
\end{enumerate}

\begin{enumerate}
\def\labelenumi{¶ \arabic{enumi})}
\setcounter{enumi}{11}
\tightlist
\item
  设 H 为可分 Hilbert 空间，\((f_n)_{n \geqslant 1}\) 为 H 的一个标准正交基。按下述方式定义诸数 \((\alpha_m)_{m \geqslant 1}\)：若 \(\alpha_m = e^{-k}\) 是 \(m\) 与一个奇数之积，则令 \(2^k\)。
\end{enumerate}

\begin{enumerate}
\def\labelenumi{\alph{enumi})}
\tightlist
\item
  设 \(x \in \mathcal{L}(\mathrm{H})\) 由 \(x(f_m) = \alpha_m f_{m+1}\)（对所有 \(m \geqslant 1\)）定义。证明 \(x\) 不是拟幂零元。（注意
\end{enumerate}

\[
\left\| x ^ {n} \right\| = \sup _ {m} \bigl (\alpha_ {m} \alpha_ {m + 1} \dots \alpha_ {m + n - 1} \bigr),
\]

并估计 \(\alpha_{1}\alpha_{2}\ldots\alpha_{2^{t}-1}\)。）

\begin{enumerate}
\def\labelenumi{\alph{enumi})}
\setcounter{enumi}{1}
\tightlist
\item
  设 \(x_{k} \in \mathcal{L}(\mathrm{H})\) 由下式定义：当 m 是 \(x_{k}(f_{m}) = 0\) 与一个奇数之积时 \(2^{k}\)，否则 \(x_{k}(f_{m}) = \alpha_{m} f_{m+1}\)。证明 \(x_{k}\) 是幂零元，且当 k 趋于 \((x_{k})\) 时 \(+\infty\) 收敛于 x。
\end{enumerate}

\begin{enumerate}
\def\labelenumi{\arabic{enumi})}
\setcounter{enumi}{12}
\item
  设 A 为幺 Banach 代数。设 \(x \in A\) 满足 \(\|x - 1\| < 1\)。则存在 \(y \in A\) 使 \(y^{2} = x\)。
\item
  设 A 为复幺赋范代数。若对 A 的每个可逆元 x 都有 \(\|x^{-1}\| = \|x\|^{-1}\)，则 A = C · 1。（可归结为 A 完备的情形。设 \(x \in A\) 与 C · 1 的距离为 \(\alpha > 0\)。若 \(\lambda_{0} \in C\) 使 \(x - \lambda_{0}\) 可逆，则当 \(x - \lambda\) 时 \(|\lambda - \lambda_{0}| < \alpha\) 也必可逆。由此推出 x 的谱将是空集。）
\item
  设 A 为复幺赋范代数，其中仅有的左拓扑零因子是 0，仅有的右拓扑零因子也是 0。则 \(A = C \cdot 1\)。（若 \(x \in A\)，且 \(\lambda\) 属于 \(\mathrm{Sp}_{\widehat{\mathrm{A}}}(\boldsymbol{x})\) 的边界，则 \(x - \lambda\) 是 \(\widehat{A}\) 中的拓扑零因子，因而是 A 中的拓扑零因子，于是 \(x = \lambda\)。）
\item
  设 A 为复幺 Banach 代数。设 \(a\) 为 A 的一个元素。我们把 a 的奇异谱称为满足下列条件的 \(\lambda \in \mathbf{C}\) 所成的集合：\(a - \lambda e\) 在 A 中既是左拓扑零因子又是右拓扑零因子。
\end{enumerate}

\begin{enumerate}
\def\labelenumi{\alph{enumi})}
\item
  证明 a 的奇异谱包含于 a 的谱中，并且包含 a 的谱的边界。
\item
  设 \(f \in C[X]\) 为多项式。a 的奇异谱在 f 下的像就是 \(f(a)\) 的奇异谱。
\end{enumerate}

\begin{enumerate}
\def\labelenumi{\arabic{enumi})}
\setcounter{enumi}{16}
\tightlist
\item
  对每个 \(x \in A\)，我们令：
\end{enumerate}

\[
\lambda (x) = \inf _ {y \neq 0} \frac {\| x y \|}{\| y \|} \quad \lambda^ {\prime} (x) = \inf _ {y \neq 0} \frac {\| y x \|}{\| y \|}.
\]

\begin{enumerate}
\def\labelenumi{\alph{enumi})}
\tightlist
\item
我们有
\end{enumerate}

\[
| \lambda (x) - \lambda (y) | \leqslant \| x - y \|, \quad | \lambda^ {\prime} (x) - \lambda^ {\prime} (y) | \leqslant \| x - y \|,
\]

\[
\lambda (x) \lambda (y) \leqslant \lambda (x y) \leqslant \| x \| \lambda (y), \quad \lambda^ {\prime} (x) \lambda^ {\prime} (y) \leqslant \lambda^ {\prime} (x y) \leqslant \lambda^ {\prime} (x) \| y \|.
\]

\begin{enumerate}
\def\labelenumi{\alph{enumi})}
\setcounter{enumi}{1}
\tightlist
\item
  A 中的左（相应地，右）拓扑零因子所成的集合是闭集。
\end{enumerate}

\begin{enumerate}
\def\labelenumi{\arabic{enumi})}
\setcounter{enumi}{17}
\item
  设 A 为幺 Banach 代数。由既不可逆又不是拓扑零因子的 \(x\) 所成的集合在 A 中是开集。
\item
  设 X 为复 Banach 空间，\(x \in A = \mathcal{L}(X)\)。
\end{enumerate}

\begin{enumerate}
\def\labelenumi{\alph{enumi})}
\tightlist
\item
  若 x 不可逆，则 x 是左拓扑零因子或右拓扑零因子。（依次考虑下列各种情形：）
\end{enumerate}

\(1^{\circ}\) ) \(x\) 不是单射；

\(2^{\circ})\overline{x(\mathrm{X})}\neq \mathrm{X};\)

\(3^{\circ}) x \text{ 是单射的， } x(\mathrm{X}) \text{ 在 X 中稠密且不同于 X。}\)

\begin{enumerate}
\def\labelenumi{\alph{enumi})}
\setcounter{enumi}{1}
\tightlist
\item
  若 x 把 X 双连续地映到 X 的一个异于 X 的闭向量子空间上，或者 x 非单射而其像为 X，则 x 位于 A 的不可逆元素所成集合的内部。
\end{enumerate}

\begin{enumerate}
\def\labelenumi{\arabic{enumi})}
\setcounter{enumi}{19}
\item
  在 I，第 20 页，例 9 的代数 A 中，恒等函数 z 既不可逆，也不是拓扑零因子。
\item
  设 A 为幺赋范代数。设 B 为赋范代数 \(\mathcal{L}(\mathrm{A})\)。于是 A 在态射 \(x \mapsto \gamma_x\) 下的像是 B 的一个全子代数。因此，若 \(x \in \mathrm{A}\)，则 \(\mathrm{Sp}_{\mathrm{A}}(x) = \mathrm{Sp}_{\mathrm{B}}(\gamma_x)\)。
\end{enumerate}

\begin{enumerate}
\def\labelenumi{¶ \arabic{enumi})}
\setcounter{enumi}{21}
\tightlist
\item
  设 A 为 Banach 代数。
\end{enumerate}

\begin{enumerate}
\def\labelenumi{\alph{enumi})}
\tightlist
\item
  设 S 为由 A 的元素组成的有界序列全体，赋以范数 \(\|(\boldsymbol{x}_{n})\| = \sup \|x_{n}\|\)。设 R 为 \((x_{n}) \in \mathrm{S}\) 中使得 \(\|x_{n}\|\) 趋于 0 的元素所成的集合。于是 S 是 Banach 代数，而 R 是 S 的一个闭双边理想。
\end{enumerate}

设 \(S' = \text{S/R}\)，并设 \(\varphi: \text{S} \to \text{S}'\) 为典范同态。

\begin{enumerate}
\def\labelenumi{\alph{enumi})}
\setcounter{enumi}{1}
\item
  映射 \(\theta\)（由 \(\theta(x) = \varphi(x, x, x, \ldots)\) 定义，从 A 到 S' 中）是 A 到 S' 的一个子代数上的等距同构。S' 中的每个左拓扑零因子都是左零因子。为使 \(x \in A\) 是左拓扑零因子，充要条件是 \(\theta(x)\) 是 S' 中的左零因子。
\item
  假设 A 有单位元。证明：存在一个 Banach 代数 B 以及 A 到 B 的一个子代数上的等距同构，使得 A 的元素不可逆当且仅当它在 B 中的像是左零因子或右零因子。（利用 a) 以及习题 19 与 21。）
\end{enumerate}

\begin{enumerate}
\def\labelenumi{\arabic{enumi})}
\setcounter{enumi}{22}
\tightlist
\item
  设 A 为赋范代数，R 为其根。
\end{enumerate}

\begin{enumerate}
\def\labelenumi{\alph{enumi})}
\item
  若 \(x \in R\)，则 x 是拟幂零元。（我们有 \(\mathrm{Sp}_{\mathrm{A}}^{\prime}(x) = \{0\}\)，从而更有 \(\mathrm{Sp}_{\widehat{\mathrm{A}}}^{\prime}(x) = \{0\}\)。）
\item
  假设 A 是 Banach 代数。设 I 为 A 的一个左理想，其所有元素都是拟幂零元。则 \(I \subset R\)（可以假设 A 有单位元。若 \(x \in R\) 且 \(a \in A\)，则有 \(\mathrm{Sp}_{\mathrm{A}}(ax) = \{0\}\)，因而 1 - ax 可逆；从而 \(x \in R\)。
\end{enumerate}

\begin{enumerate}
\def\labelenumi{\arabic{enumi})}
\setcounter{enumi}{23}
\item
  设 A 为复 Banach 代数，B 为无根的复代数，\(\varphi\) 为从 A 到 B 上的态射。\(\varphi\) 的核是闭的。
\item
  设 A 为复 Banach 代数，\(\pi\) 为 A 在复向量空间 X 中的一个非零不可约表示。设 \(\xi_{0}\) 为 X 的非零元素。
\end{enumerate}

\begin{enumerate}
\def\labelenumi{\alph{enumi})}
\item
  \(\xi_0\) 的零化子 I 是 A 的一个正则极大左理想；它是闭的。
\item
  映射 \(x \mapsto x\xi_{0}\) 通过过渡到商，定义了 A-模 A/I 到 A-模 X 上的一个同构 \(\varphi\)。
\item
  若经 \(\varphi\) 把 Banach 空间 A/I 的范数转移过来，则 X 成为 Banach 空间，且对所有 \(\| \pi (x)\| \leqslant \| x\|\) 有 \(x\in \mathrm{A}\)。
\item
  若 A 是本原的（即 \(\{0\}\) 是本原理想，参见 I，第 11 页的定义，另参见 A，VIII，第 433 页，习题 6），则依照 A 有或没有单位元而分别有 \(A = \{0\}\) 或 \(C \cdot 1\)。（利用前面的结果以及 I，第 26 页的推论 4。）
\end{enumerate}

\begin{enumerate}
\def\labelenumi{¶ \arabic{enumi})}
\setcounter{enumi}{25}
\tightlist
\item
  设 A 为 Banach 代数，R 为其根。
\end{enumerate}

\begin{enumerate}
\def\labelenumi{\alph{enumi})}
\tightlist
\item
  证明：若 \(r \in R\)，则级数
\end{enumerate}

\[
- \frac {1}{2} \sum_ {k = 1} ^ {\infty} \binom{1 / 2}{k} (- 4 r) ^ {k}
\]

收敛于满足 \(x \in R\) 的元素 \(x^{2} - x + r = 0\)。

\begin{enumerate}
\def\labelenumi{\alph{enumi})}
\setcounter{enumi}{1}
\tightlist
\item
  设 u 为 A 的一个元素，它模 R 的类是幂等的。则存在 A 中与 u 模 R 同余的一个幂等元。（可以假设 A 有单位元。设 \(q = u - u^{2} \in R\)，并设 \(x \in R\) 为借助 a) 构造的方程 \(x^{2} - x - q(1 - 4q)^{-1} = 0\) 的一个解。那么 \(u - (2u - 1)x\) 即为所求。）
\end{enumerate}

\begin{enumerate}
\def\labelenumi{\arabic{enumi})}
\setcounter{enumi}{26}
\item
  设 A 为复幺 Banach 代数，\(x \in A\)，并设 \(\zeta\) 为 \(\mathrm{Sp}_{\mathrm{A}}(x)\) 的一个边界点。x 的预解式不能延拓为在 \(\zeta\) 处连续的函数。
\item
  设 A 为复幺 Banach 代数，\(\Delta\) 为 C 的开子集，\(\lambda \mapsto \mathrm{R}(\lambda)\) 为从 \(\Delta\) 到 A 中的映射，使得
\end{enumerate}

\[
\mathrm{R} (\lambda) - \mathrm{R} (\mu) = - (\lambda - \mu) \mathrm{R} (\lambda) \mathrm{R} (\mu)\tag{5}
\]

对所有 \(\lambda,\mu\in\Delta\) 成立。

\begin{enumerate}
\def\labelenumi{\alph{enumi})}
\tightlist
\item
  设 \(\lambda_0\in \Delta\)。对每个满足 \(\lambda \in \Delta\) 的 \(|\lambda -\lambda_0|\cdot \| \mathrm{R}(\lambda)\| <  1\)，我们有：
\end{enumerate}

\[
\mathrm{R} (\lambda) = \sum_ {n = 0} ^ {\infty} (\lambda_ {0} - \lambda) ^ {n} \mathrm{R} (\lambda_ {0}) ^ {n + 1}.
\]

\begin{enumerate}
\def\labelenumi{\alph{enumi})}
\setcounter{enumi}{1}
\tightlist
\item
  对每个 \(\lambda \in \Delta\) 和每个整数 \(n\geqslant 0\)，有
\end{enumerate}

\[
\frac {\partial^ {n}}{\partial \lambda^ {n}} \mathrm{R} (\lambda) = (- 1) ^ {n} n! \mathrm{R} (\lambda) ^ {n + 1}.
\]

\begin{enumerate}
\def\labelenumi{\alph{enumi})}
\setcounter{enumi}{2}
\tightlist
\item
  若 R 在无穷远处全纯，则存在 \(z, j, x \in \mathrm{A}\) 使得 \(z^2 = 0\)、\(j^2 = j\)、\(zj = jz = 0\)、\(x \in j\mathrm{Aj}\)，且当 \(\mathrm{R}(\lambda) = z + j\mathrm{R}(\lambda, x)\) 充分大时 \(|\lambda|\)。（当 \(\mathrm{R}(\lambda) = \sum_{n=0}^{\infty} c_n \lambda^{-n}\) 时写出 \(|\lambda| > \lambda_0\)，并表出 R 满足 (5)。可令 \(c_0 = z\)、\(c_1 = j\)、\(c_2 = x\)。）
\end{enumerate}

\begin{enumerate}
\def\labelenumi{\alph{enumi})}
\setcounter{enumi}{3}
\tightlist
\item
  为存在 \(a \in A\) 使 \(R\) 是 \(\Delta\) 的预解式 \(\lambda \mapsto R(a, \lambda)\) 在 \(a\) 上的限制，充要条件是存在 \(\lambda_0 \in \Delta\) 使 \(R(\lambda_0)\) 可逆。
\end{enumerate}

\begin{enumerate}
\def\labelenumi{\alph{enumi})}
\setcounter{enumi}{4}
\tightlist
\item
  对每个 \(x \in A\)，函数 \(\lambda \mapsto \sum_{n=0}^{\infty} (\lambda_0 - \lambda)^n x^{n+1}\) 在条件 \(|\lambda - \lambda_0| \cdot \|x\| < 1\) 下满足 (5)。
\end{enumerate}

\begin{enumerate}
\def\labelenumi{\arabic{enumi})}
\setcounter{enumi}{28}
\tightlist
\item
  \begin{enumerate}
  \def\labelenumii{\alph{enumii})}
  \tightlist
  \item
    设 E 为复 Banach 空间，\(f: \mathbf{C} \to \mathbf{E}\) 为整函数，\(\theta \mapsto \mathrm{P}(\theta) = \sum_{\mu=k}^{m} c_{\mu} e^{\mu i\theta}\) 为三角多项式（\(c_k \neq 0\)）。假设当 \(|\mathrm{P}(\theta)| \cdot \|f(re^{i\theta})\| \leqslant \mathrm{Mr}^{\alpha}\) 时 \(r \geqslant r_0\)（M、\(r_0\)、\(\alpha\) 为常数，均 \(\geqslant 0\)）。则 \(f\) 是次数 \(\leqslant \alpha\) 的多项式。（设 \(f(\zeta) = \sum_{\nu=0}^{\infty} a_{\nu} \zeta^{\nu}\)，并令
  \end{enumerate}
\end{enumerate}

\[
\mathrm{J} (r, p) = r ^ {- p} \int_ {0} ^ {2 \pi} \mathrm{P} (\theta) f (r e ^ {i \theta}) e ^ {- (k + p) i \theta} d \theta .
\]

对 \(\mathrm{J}(r,p)\)，用两种不同的方式计算当 \(r\to+\infty\) 时 \(p>\alpha.\) 的极限。

\begin{enumerate}
\def\labelenumi{\alph{enumi})}
\setcounter{enumi}{1}
\tightlist
\item
  沿用 \(a\) 的记号，若 \(\alpha\) 是整数，且对每个 \(\theta \in \mathbf{R}\) 有
\end{enumerate}

\[
\lim _ {r \to + \infty} r ^ {- \alpha} | \mathrm{P} (\theta) | \| f (r e ^ {i \theta}) = 0,
\]

则 \(f\) 是次数 \(< \alpha\) 的多项式。

\begin{enumerate}
\def\labelenumi{\alph{enumi})}
\setcounter{enumi}{2}
\tightlist
\item
  设 A 为复幺 Banach 代数，x 为 A 的一个拟幂零元。假设对某个整数 n \textgreater{} 0，且对每个 \(\theta \in R\)，有
\end{enumerate}

\[
\lim _ {r \to + \infty} r ^ {- n} | \cos^ {n} \theta | \cdot \| (1 - r e ^ {i \theta} x) ^ {- 1} \| = 0.
\]

则 \(x^n = 0\)。（利用 \(b\)。）

\begin{enumerate}
\def\labelenumi{¶ \arabic{enumi})}
\setcounter{enumi}{29}
\tightlist
\item
  a) 设 E 为复 Banach 空间，\(x: \mathbf{C}=\{1\}\to\mathrm{E}\) 为 \(1/(\zeta-1)\) 的整函数。用 \(x(\zeta)=\sum_{n=0}^{\infty}a_{n}\zeta^{n}\)、\(x(\zeta)=\sum_{n=0}^{\infty}b_{n}\zeta^{-n}\) 分别表示 \(x\) 在 \(|\zeta|<1\)、\(|\zeta|>1\) 中的展开式。假设存在 \(\alpha>0\) 使得 \(\|a_{n}\|=o(n^{\alpha})\)、\(\|b_{n}\|=o(n^{\alpha})\)。则 \(x(\zeta)\) 是 \(1 / (\zeta - 1)\) 的次数 \(< \alpha + 1\) 的多项式。（设 \(\varepsilon > 0\)。我们有：当 \(\|x(\zeta)\| \leqslant \varepsilon(1 - r)^{-1 - \alpha}\) 时 \(r_{\varepsilon} \leqslant r = |\zeta| < 1\)；当 \(\|x(\zeta)\| \leqslant \varepsilon(1 - r^{-1})^{-1 - \alpha}\) 时 \(1 < r \leqslant r_{\varepsilon}'\)。令 \(1 - \zeta = (\omega + \frac{1}{2})^{-1}\)、\(\omega = \mathrm{Re}^{i\theta}\)，证明：对 \(0 \leqslant r < 1\) 与 \(\mathrm{R} \geqslant 1\)，有
\end{enumerate}

\[
(1 - r) ^ {- 1} \leqslant \frac {9}{4} \frac {\mathrm{R}}{\cos \theta}
\]

并且，对 1 \textless{} r 与 R ≥ 1，

\[
(1 - r ^ {- 1}) ^ {- 1} \leqslant \frac {9}{4} \frac {\mathrm{R}}{\cos \theta}.
\]

令 \(y(\omega) = x(\zeta)\)，并证明当 \(R^{-\alpha - 1}|\cos \theta|^{\alpha + 1}\| y(Re^{i\theta})\|\) 趋于 \(R\) 时 \(+\infty\) 关于 \(\theta\) 一致地趋于 0。然后应用习题 29 的 b)。

设 A 为复幺 Banach 代数，\(q \in A\) 为拟幂零元，且 \(x = 1 + q\)。

\begin{enumerate}
\def\labelenumi{\alph{enumi})}
\setcounter{enumi}{1}
\item
  为使 \(q^{N}=0\)，充要条件是当 n 趋于 \(\|x^{\pm n}\|=o(n^{N})\) 时 \(+\infty\)。（为看出条件是充分的，应用 a) 证明 \(\mathrm{R}(\lambda,x)\) 是 \(1/(\lambda-1)\) 的次数 \(\leqslant N\) 的多项式。另一方面，\(\mathrm{R}(\lambda,x)=\sum_{n=0}^{\infty}q^{n}(\lambda-1)^{-n-1}\)。）
\item
  设 R 为 A 的根，G 为 A 的可逆元群的一个子群。若对每个 \(x \in G\)，当 \(\|x^{n}\| = o(n)\) 时有 \(\|x^{-n}\| = o(n)\)、\(n \to +\infty\)，则典范映射 \(A \to A/R\) 在 G 上的限制是单射。
\end{enumerate}

\begin{enumerate}
\def\labelenumi{\arabic{enumi})}
\setcounter{enumi}{30}
\tightlist
\item
  \begin{enumerate}
  \def\labelenumii{\alph{enumii})}
  \tightlist
  \item
    设 A 为赋范代数，x、y 为 A 中满足 xy - yx = y 的两个元素。证明：当 \(y^{n} = 0\) 时 \(n > 2\|x\|\)。（利用公式 \(xy^{n} - y^{n}x = ny^{n}\)。）
  \end{enumerate}
\end{enumerate}

\begin{enumerate}
\def\labelenumi{\alph{enumi})}
\setcounter{enumi}{1}
\item
  设 L 为 C 上的有限维 Lie 代数。证明：若 L 不是幂零的，则它含有两个非零元素 x、y，使得 \([x, y] = y\)。（利用存在 \(x \in L\) 使 \(\operatorname{ad}(x)\) 不是幂零的这一事实。）
\item
  设 L 为实或复的有限维 Lie 代数，且 L 的包络代数具有一个与其代数结构相容的范数。证明 L 是幂零的。（利用 a) 与 b)。）
\end{enumerate}

\begin{enumerate}
\def\labelenumi{\alph{enumi})}
\setcounter{enumi}{3}
\tightlist
\item
  设 V 为实或复的有限维向量空间，\(\mathfrak{L}\) 为 \(\mathfrak{gl}(\mathrm{V})\) 的一个由幂零自同态组成的 Lie 子代数，\(x_{1},\ldots ,x_{n}\) 为 \(\mathfrak{L}\) 中的元素。设 \(\alpha >0\)。证明：V 上存在一个 Hilbert 空间结构，使得对 \(\| x_i\| \leqslant \alpha\) 有 \(1\leqslant i\leqslant n\)。
\end{enumerate}

\begin{enumerate}
\def\labelenumi{\alph{enumi})}
\setcounter{enumi}{4}
\tightlist
\item
  设 \(\mathfrak{L}\) 为实或复的有限维幂零 Lie 代数，U 为其包络代数。证明：U 上存在一个与其代数结构相容的范数。（利用 LIE，I，§3，习题 5 与 §7，习题 3 的方法，证明存在一族表示 \((\pi_{\lambda})\)，它们是 \(\mathfrak{L}\) 在有限维空间 \(V_{\lambda}\) 中的表示，具有下列性质：
\end{enumerate}

对所有 \(u \in U\)，存在 \(\lambda\) 使 \(\pi_{\lambda}(u) \neq 0\)；并且对每个 \(\pi_{\lambda}(\mathfrak{L})\)，\(\lambda\) 都由幂零自同态组成。给每个 \(V_{\lambda}\) 赋以由 d) 提供的 Hilbert 结构，便在 Hilbert 空间 V（诸 \(V_{\lambda}\) 的 Hilbert 和）中得到 L 的一个由连续自同态给出的表示 \(\pi\)，从而得到一个单射态射 \(\varphi\)，它把 U 映入 \(\mathcal{L}(V)\)。令 \(\|u\|_{U} = \|\varphi(u)\|\)（对所有 \(u \in U\)）。

\begin{enumerate}
\def\labelenumi{\alph{enumi})}
\setcounter{enumi}{5}
\tightlist
\item
  设 A 为复幺赋范代数，x、y 为 A 的元素。若 \(A \neq \{0\}\)，则 xy - yx ≠ 1。（第一个证明：我们有 \(\mathrm{Sp}_{\mathrm{A}}^{\prime}(xy) = \mathrm{Sp}_{\mathrm{A}}^{\prime}(yx)\)；若 xy = yx + 1，则 \(\mathrm{Sp}_{\mathrm{A}}(xy)\) 可由 \(\mathrm{Sp}_{\mathrm{A}}(yx)\) 经平移 \(z \mapsto z + 1\) 得到；由此推出 \(\mathrm{Sp}_{\mathrm{A}}(xy)\) 无界，而这是荒谬的。第二个证明：若 \([x, y] = 1\)，则有 \([xy, y] = y\)，从而由 a) 知对充分大的 n 有 \(y^{n} = 0\)；再由 \([x, y^{p}] = py^{p-1}\) 推出 y = 0。）
\end{enumerate}

\begin{enumerate}
\def\labelenumi{\arabic{enumi})}
\setcounter{enumi}{31}
\item
  设 A 为赋范代数，\(x \in A\)。若当 n 趋于 \(x^{n}\) 时 \(+\infty\) 趋于 0，则称 x 是拓扑幂零的。证明：x 是拓扑幂零的当且仅当 \(\varrho(x) < 1\)。
\item
  设 A 为 C 上的交换幺 Banach 代数。假设 A 的每个闭理想 I 都是有限生成的（即存在 \(n \in N\) 与 \(x_{1}, \ldots, x_{n} \in A\)，使得 \(I = Ax_{1} + \cdots + Ax_{n}\)）。
\end{enumerate}

\begin{enumerate}
\def\labelenumi{\alph{enumi})}
\tightlist
\item
  A 的每个理想 I 都是闭的。（写 \(\bar{\mathrm{I}} = \mathrm{A}x_{1} + \dots +\mathrm{A}x_{n}\)，其中 \(x_{1},\ldots ,x_{n}\) 属于 A。设 \((x_{ip})_{p\geqslant 1}\) 是 I 中趋于 \(x_{i}\) 的元素序列。对每个 \((y_{1},\ldots ,y_{n})\in \mathrm{A}^{n}\) 与每个 \(p\geqslant 1\)，令
\end{enumerate}

\[
\begin{array}{c} \varphi (y _ {1}, \ldots , y _ {n}) = x _ {1} y _ {1} + \dots + x _ {n} y _ {n} \in \bar {\mathrm{I}} \\ \varphi_ {p} (y _ {1}, \ldots , y _ {n}) = x _ {i p} y _ {1} + \dots + x _ {n p} y _ {n} \in \mathrm{I}. \end{array}
\]

我们有 \(\varphi\in\mathcal{L}(A^{n},\bar{I})\)、\(\varphi_{p}\in\mathcal{L}(A^{n},\bar{I})\)，且 \(\|\varphi-\varphi_{p}\|\) 随 \(p\) 趋于 \(+\infty\) 而趋于 0，并且 \(\varphi\) 是满射的。借助 \(\varphi_{p}\) 当 \(p\) 充分大时的满射性（考虑 \(\varphi\) 与 \(\varphi_{p}\) 的转置），推出 I = \(\bar{I}\)。

\begin{enumerate}
\def\labelenumi{\alph{enumi})}
\setcounter{enumi}{1}
\item
  若 A 是整环，则 \(A = C \cdot 1\)。（利用 a) 与习题 15。）
\item
  若 A 中唯一的幂零元是 0，则存在整数 \(n \geqslant 0\) 使 A 同构于 \(\mathbf{C}^n\)。（由于 A 是 Noether 的，\(\{0\}\) 是素理想 \(\mathfrak{p}_1, \ldots, \mathfrak{p}_m\) 的交；由 \(a\)) 知它们都是闭的；而代数 \(\mathrm{A}/\mathfrak{p}_i\) 同构于 \(\mathbf{C}\)（据 \(b\)）。）
\end{enumerate}

\begin{enumerate}
\def\labelenumi{\alph{enumi})}
\setcounter{enumi}{3}
\tightlist
\item
  我们有 \(\dim_{\mathbf{C}}\mathrm{A} < +\infty\)。（设 N 为 A 的幂零元全体，它是 A 的一个（闭）理想。于是由 \(\dim (\mathrm{A / N}) < +\infty\)) 有 \(c\)。由于 N 是有限生成理想，故对充分大的 \(\mathrm{N}^n = 0\) 有 \(n\)。最后，每个 \(\mathrm{N}^i /\mathrm{N}^{i + 1}\) 都是 A/N 上的有限生成模，故 \(\dim (\mathrm{N}^i /\mathrm{N}^{i + 1}) < +\infty\)。）
\end{enumerate}

\begin{enumerate}
\def\labelenumi{\arabic{enumi})}
\setcounter{enumi}{33}
\tightlist
\item
  设 A 为 C 上的幺代数，带有一个分离的局部凸拓扑，使得 A 中的乘法分别连续。
\end{enumerate}

。称元素 \(a \in A\) 为正则的，如果存在 \(r \geqslant 0\)，使得 \(a - \lambda\) 对 \(|\lambda| \geqslant r\) 可逆，且使得 \((a - \lambda)^{-1}\)（对 \(|\lambda| \geqslant r\)）所成的集合有界。

\begin{enumerate}
\def\labelenumi{\alph{enumi})}
\item
  若 a 正则，则 \(\lambda(a-\lambda)^{-1}\) 对 \(|\lambda|\geqslant r\) 有界。
\item
  若映射 \(x \mapsto x^{-1}\) 在 1 的一个邻域中有定义且在点 1 处连续，则 A 的每个元素都是正则的。
\item
  对每个 \(a \in A\)，设 \(U_{a}\) 为使 \(\lambda \in C\)，且 \((a - \mu)^{-1}\)（对 \(\mu\) 取遍 \(\lambda\) 的一个邻域）存在且有界的元素所成的集合。设 \(S_{a} = C - U_{a}\)。则 \(S_{a}\) 是闭的；且若 a 正则，则 \(S_{a}\) 是紧的且非空（假定 \(A \neq \{0\}\)）。
\item
  设 \(a \in A\)。函数 \(\lambda \mapsto (a - \lambda)^{-1}\) 在 \(U_{a}\) 中全纯。
\item
  设 \(a \in A\)。为使 \(\lambda \in U_{a}\)，充要条件是 \(a - \lambda\) 有正则逆。
\item
  若 A 的每个元素都正则，且 A 是域，则 \(A = C \cdot 1\)。
\end{enumerate}

\begin{enumerate}
\def\labelenumi{\arabic{enumi})}
\setcounter{enumi}{34}
\tightlist
\item
  设 A 为 R 上的一个代数且是域，带有一个分离的局部凸拓扑，使得 A 中的乘法连续，且映射 \(x \mapsto x^{-1}\) 在 1 处连续。则 A 与 R、与 C 或与四元数域 H 是 R-同构的。（若 A 交换且存在 \(u \in A\) 满足 \(u^{2} = -1\)，则给 A 一个 C 上的代数结构，并应用习题 34 的 f)。）若 A 交换且不存在元素 \(u \in A\) 满足 \(u^{2} = -1\)，则对域 \(A \otimes_{R} C\) 应用习题 34 的 f)。若 A 非交换，则按 AC，VI，§6，第 4 节，定理 1 的第三种情形论证。
\end{enumerate}

\begin{enumerate}
\def\labelenumi{¶ \arabic{enumi})}
\setcounter{enumi}{35}
\tightlist
\item
  设 A 为由复系数单变量有理函数到 {[}0,1{]} 上的限制所组成的 C 上代数。这个代数是一个域。
\end{enumerate}

\begin{enumerate}
\def\labelenumi{\alph{enumi})}
\item
  给 A 赋以依测度收敛的拓扑（INT，IV，§5，第 11 节）。则映射 \((x,y)\mapsto x+y\) 与 \((x,y)\mapsto xy\)（从 A × A 到 A 中）连续，且映射 \(x\mapsto x^{-1}\)（从 A* 到 A 中）连续。A 的拓扑不是局部凸的。
\item
  对 \(n \in N^{*}\)，定义序列 \((w_{r,n})_{r \in \mathbf{Z}}\) 如下：若 r \(w_{r,n} = (-r + 1)^{n(-r+1)}\)\textless{} 0，\(w_{0,n} = 1\)，且若 r \(w_{r,n} = (r + 1)^{-(r+1)/n}\)\textgreater{} 0。若 \(f \in A\)，令 \(p_n(f) = \sum_r w_{r,n} |a_r| < +\infty\)，其中 \(\sum a_r t^r\) 是 \(f(t)\) 在 0 处的 Laurent 展开式。于是 \(p_n\) 是一些半范数，它们在 A 上定义了一个可度量的局部凸拓扑。映射 \((x, y) \mapsto xy\)（从 A × A 到 A 中）连续。映射 \(x \mapsto x^{-1}\)（从 A* 到 A 中）不连续。代数 A 对此拓扑不完备。
\end{enumerate}

\begin{enumerate}
\def\labelenumi{\arabic{enumi})}
\setcounter{enumi}{36}
\tightlist
\item
  \begin{enumerate}
  \def\labelenumii{\alph{enumii})}
  \tightlist
  \item
    设 A 为 C 上带有局部凸拓扑的一个代数。下列条件等价：
  \end{enumerate}
\end{enumerate}

\begin{enumerate}
\def\labelenumi{(\roman{enumi})}
\item
  存在 0 的凸平衡邻域的基本组 \((\mathrm{U}_{i})\)，使得对一切 i 有 \(U_{i}U_{i} \subset U_{i}\)；
\item
  A 同构于一族赋范代数的乘积的一个子代数。
\item
  A 的拓扑可由一族半范数 \(p_{i}\) 定义，\(p_{i}(xy) \leqslant p_{i}(x)p_{i}(y)\) 对一切 \(x, y \in A\) 成立。
\end{enumerate}

若 A 满足这些条件，则称 A 是局部 m-凸的。

\begin{enumerate}
\def\labelenumi{\alph{enumi})}
\setcounter{enumi}{1}
\item
  设 A 为局部 m-凸代数。映射 \((x,y)\mapsto xy\)（从 A × A 到 A 中）连续。若 A 有单位元，且用 G 表示 A 的可逆元集合，则映射 \(x\mapsto x^{-1}\)（从 G 到 A 中）连续。
\item
  设 A 为幺局部 m-凸代数。A 的每个元素的谱非空。若 A 是域，则 \(A = C \cdot 1\)。
\item
  设 A 为完备的局部 m-凸代数。存在一个有向集 I、一族 Banach 代数 \((\mathrm{A}_{i})_{i\in\mathrm{I}}\)、以及对 \(\pi_{ij}\colon A_{j}\mapsto A_{i}\) 的连续态射 \(i\leqslant j\)，满足 \(\pi_{ij}\circ\pi_{jk}=\pi_{ik}\)，使得 A 同构于 \(\prod_{i}A_{i}\) 的由那些 \((x_{i})\)（它们满足 \(\pi_{ij}(x_{j})=x_{i}\)，对 \(i\leqslant j\)）组成的子代数。
\item
  代数 \(\mathcal{C}_{\mathbf{C}}(\mathbf{R})\)，赋以紧收敛拓扑，是局部 m-凸的。\(\mathcal{C}_{\mathbf{C}}(\mathbf{R})\) 的可逆元集合在 \(\mathcal{C}_{\mathbf{C}}(\mathbf{R})\) 中稠密；它不是开集。
\item
  设 D 为 C 的开子集，A 为 D 上复值全纯函数所成的代数，赋以紧收敛拓扑。则 A 是局部 m-凸的。A 的可逆元集合是闭的。
\item
  {[}0,1{]} 上无限次可微的复值函数所成的代数，赋以各阶导数一致收敛的拓扑，是局部 m-凸的。
\item
  设 A 为 {[}0,1{]} 上满足：对每个 p \(f \in \mathrm{L}^{p}([0,1])\)\textgreater{ 1 有 } 的复函数（的类）所成的代数。赋以由半范数族 \(f \mapsto \|f\|_{p}\)（p \textgreater{} 1）定义的拓扑，A 不是局部 m-凸的。映射 \((x,y) \mapsto xy\)（从 A × A 到 A 中）连续。
\end{enumerate}

